% Teil II: der durchgerechnete Weg.

\subsection*{Gegeben und gesucht}
(a) Programm so ändern, dass $0+1+\dots+29$ summiert wird; (b) Programm
für $n=1,2,3,4,\dots$ auswerten und das Muster beschreiben; (c)
$\sum_{k=0}^{n}k=n(n+1)/2$ für alle $n\ge1$ beweisen.

\subsection*{a) Programmänderung}
\emph{Vorher:} \texttt{sum(range(21))} (summiert $0,\dots,20$).
Für $0,\dots,29$:
\begin{verbatim}
sum(range(30))
\end{verbatim}

\subsection*{b) Auswertung von \texttt{sum(range(n+1))-n*(n+1)/2}}
\begin{center}
\begin{tabular}{c|c|c|c}
$n$ & \texttt{sum(range(n+1))} & $n(n+1)/2$ & Differenz\\
\hline
1 & 1 & 1 & 0\\
2 & 3 & 3 & 0\\
3 & 6 & 6 & 0\\
4 & 10 & 10 & 0
\end{tabular}
\end{center}
Das Ergebnis ist für jedes $n$ \emph{stets $0$}, denn beide Terme stimmen
überein.

\subsection*{c) Beweis per vollständiger Induktion}
\emph{Induktionsanfang} ($n=1$): $\sum_{k=0}^{1}k=0+1=1=\frac{1\cdot2}2$. \\
\emph{Induktionsschritt} ($n\to n+1$), mit der Voraussetzung
$\sum_{k=0}^{n}k=\frac{n(n+1)}2$:
\begin{align}
  \sum_{k=0}^{n+1}k &=\sum_{k=0}^{n}k+(n+1)
  =\frac{n(n+1)}2+(n+1) \\
  &=\frac{n(n+1)+2(n+1)}2=\frac{(n+1)(n+2)}2 .
\end{align}
Das ist die Behauptung für $n+1$. Aus IA und IS hält die Formel für alle
$n\ge1$.

\begin{ergebnis}
\begin{equation*}
  \sum_{k=0}^{n}k=\frac{n(n+1)}2 \quad\text{für alle } n\ge1 .
\end{equation*}
\end{ergebnis}

\begin{kontrolle}
Für $n=4$: $0+1+2+3+4=10=\frac{4\cdot5}2$; für $n=5$: $15=\frac{5\cdot6}2$.
Passt. Damit liefert das Python-Programm identisch $0$.
\end{kontrolle}

\begin{merksatz}
``Bis $n$'' heißt \texttt{range(n+1)}; die Gauß-Formel beweist man, indem
man im Induktionsschritt den letzten Summanden $(n+1)$ abspaltet und auf
den gemeinsamen Nenner bringt.
\end{merksatz}
